
\section{结论：预告后盈利桥已经闭合}

工业富联不再只是AI算力需求锚，而是本报告第七只正式估值标的。公司预计
2026H1归母净利润234--244亿元，中值239亿元、同比约+97\%；其中Q1为
105.95亿元，隐含Q2约133.05亿元、环比+25.6\%。更重要的是，公司同时披露
云服务商AI服务器收入同比增长超过230\%、800G以上数据中心交换机出货量
同比增长1.4倍，盈利增长有直接业务代理，而不是单纯依靠估值扩张。

\begin{dashboardbox}[工业富联：预告后当前价格估值]
\begin{tabularx}{\textwidth}{L{2.4cm}R{1.7cm}R{1.7cm}R{1.7cm}R{1.7cm}R{1.8cm}X}
\toprule
\textbf{现价} & \textbf{H1净利} & \textbf{2026E净利} & \textbf{EPS} &
\textbf{基准价值} & \textbf{最终目标} & \textbf{空间/动作} \\
\midrule
66.27元 & 239亿元 & 560.1亿元 & 2.82元 & 84.60元 & 82.29元 &
+24.2\%；业绩兑现后的回撤配置 \\
\bottomrule
\end{tabularx}
\end{dashboardbox}

\section{预告质量：Q2继续加速，主营证据最完整}

H1预告EPS中值约1.20元，相当于华泰预告后2026E EPS 2.82元的42.4\%。
这一进度略快于简单时间进度，但本报告没有把H1机械翻倍。华泰维持2026E
营收14,803.7亿元、归母净利润560.1亿元；华创给出更高的645.97亿元净利
预测。AStock采用华泰作为基准分母，将华创仅作为牛市盈利交叉验证。

\begin{exhibitbox}[Exhibit 11：收入代理到EPS的闭环]
\begin{tabularx}{\textwidth}{L{3.0cm}L{3.0cm}L{3.2cm}X}
\toprule
\textbf{驱动} & \textbf{已披露证据} & \textbf{盈利转换} & \textbf{仍缺字段} \\
\midrule
AI服务器 & H1云服务商收入同比+230\%以上 & 产品结构和系统价值量提升，推升收入与净利 &
客户分配、单机ASP和分部毛利率未披露 \\
高速网络 & 800G以上交换机出货同比+140\% & 高速互联占比提升改善通信业务质量 &
具体收入占比与价格未披露 \\
下一代平台 & 与大客户联合研发，H2逐步量产 & 支撑H2收入和利润继续增长 &
量产爬坡节奏与良率未披露 \\
现金转换 & Q1经营现金流250.24亿元 & 当前现金流覆盖利润，降低纯账面增长风险 &
H2营运资金仍受大客户和备货影响 \\
\bottomrule
\end{tabularx}
\sourcenote{公司2026H1预告、2026Q1财务包、华泰和华创2026-07-10预告点评。}
\end{exhibitbox}

\section{券商原文穿透与AStock冷判断}

\textbf{华泰证券}在预告后维持“买入”和93元目标，预测2026--2028营收
14,803.7/18,837.1/21,890.3亿元，归母净利润560.1/674.9/779.6亿元，
给予2026E 33倍PE。其核心假设是GB300/Rubin平台放量、单机价值量提升、
Hyperscaler资本开支延续以及液冷自供比例提高。

\textbf{华创证券}同日维持“强推”，预测2026--2028归母净利润
645.97/868.66/1,094.91亿元。归档页面没有披露点目标，因此该行只用于
盈利牛市情景，不进入Street目标权重；这避免把“强推”评级伪造成目标价。

AStock基准只给30倍PE而非华泰33倍，得到84.60元基本面价值；再按65\%
基本面、25\%市场锚72元、10\%华泰目标93元加权，最终目标82.29元。现价
对应华泰EPS约23.5倍PE，仍有空间，但已经不是低位静默布局。

\section{催化、风险与失效条件}

\begin{itemize}
  \item \textbf{升级催化}：H2下一代AI服务器进入稳定量产；AI服务器收入增速仍高于100\%；
  800G以上交换机维持增长；毛利率持续高于7\%并改善。
  \item \textbf{基准门槛}：华泰全年560.1亿元意味着H2至少贡献约321.1亿元净利；
  若H2显著低于该水平，30倍基准PE不成立。
  \item \textbf{主要风险}：北美云厂资本开支下修、GPU平台量产延期、客户集中、贸易摩擦、
  低毛利制造属性重新主导估值，以及营运资金波动。
  \item \textbf{动作纪律}：62--66元为第一观察区；72--76元进入前高压力区；未出现H2
  盈利与毛利验证前，不以93元券商目标作为无条件追涨理由。
\end{itemize}
